\section{Conclusion}

In this work, we introduced \textsf{StructSynth}, a two-stage framework that addresses the challenge of high-fidelity tabular data synthesis in low-data regimes through a \emph{discover-then-synthesize} paradigm. The core insight is to decouple dependency structure discovery from data generation: the first stage constructs a Directed Acyclic Graph (DAG) via LLM-guided breadth-first search augmented with statistical association cues, serving as a generative blueprint rather than a causal claim; the second stage leverages this blueprint to synthesize data autoregressively through topological layering, guaranteeing adherence to feature dependencies by construction.

Extensive experiments across six real-world datasets demonstrate that \textsf{StructSynth} achieves state-of-the-art downstream utility with a perfect average rank of 1.00. Notably, it improves upon the strongest baseline (CLLM) by +1.65 AUC points on average, while simultaneously achieving the best privacy preservation (avg.\ rank 1.33) and maintaining competitive statistical fidelity. Our evaluations on three benchmark DAGs with known ground-truth structures confirm that the LLM-guided discovery reliably recovers dependency structures even from as few as 20 samples. Furthermore, the consistent performance gains across seven diverse LLM architectures---spanning both open-source and proprietary models---validate the generalizability of our structural guidance paradigm.

Looking forward, several promising directions emerge from this work. First, extending the discovery stage to handle high-dimensional feature spaces ($K > 50$) through hierarchical or block-wise strategies would broaden the framework's applicability. Second, exploring more expressive structural representations beyond DAGs, such as chain graphs, could better capture bidirectional relationships. Third, distilling the discovery and generation capabilities into locally hosted open-source models would reduce API costs and improve reproducibility. Finally, applying the discover-then-synthesize paradigm to specialized objectives such as fairness-aware generation and differentially private synthesis represents a natural extension of this work.

%%
